% RSVD-CBE as the basis update of the BUG integrator.
%
% Companion to:
%   docs/current_work.md          -- the living status document
%   docs/cbe_bug.md               -- the design and its reference implementations
%   examples/cbe_walkthrough.jl   -- produces every number in Part III
%
% Compile:  pdflatex rsvd_cbe_writeup.tex   (twice, for references)
% No exotic packages: amsmath, amssymb, booktabs, listings, geometry, hyperref, xcolor.

\documentclass[11pt,a4paper]{article}

\usepackage[margin=2.4cm]{geometry}
\usepackage{amsmath,amssymb}
\usepackage{booktabs}
\usepackage{listings}
\usepackage{xcolor}
\usepackage[hidelinks]{hyperref}
\usepackage{enumitem}

\definecolor{codebg}{gray}{0.96}
\definecolor{keep}{rgb}{0.10,0.45,0.20}
\definecolor{drop}{rgb}{0.65,0.10,0.10}

\lstset{
  basicstyle=\ttfamily\footnotesize,
  backgroundcolor=\color{codebg},
  frame=single,
  framerule=0pt,
  columns=fullflexible,
  keepspaces=true,
  breaklines=true,
  postbreak=\mbox{\textcolor{drop}{$\hookrightarrow$}\space},
  showstringspaces=false,
  xleftmargin=0.4em,
  aboveskip=0.8em,
  belowskip=0.8em,
}

% Julia is not guaranteed to be predefined by `listings` (it depends on the installed
% version), and an undefined language is a hard error. Define it here so the document
% compiles against any listings release; a redefinition is harmless if one already exists.
\lstdefinelanguage{julia}{
  morekeywords={function,end,return,if,else,elseif,for,while,in,const,local,global,
                struct,mutable,using,import,do,let,begin,true,false,nothing,Int,Float64},
  sensitive=true,
  morecomment=[l]{\#},
  morestring=[b]",
}

\lstdefinestyle{jl}{
  language=julia,
  keywordstyle=\color{keep}\bfseries,
  commentstyle=\color{drop}\itshape,
}

% Notation used throughout.
\newcommand{\Hb}{\hat{H}}
\newcommand{\Theb}{\Theta}
\newcommand{\Om}{\Omega}
\newcommand{\Pperp}{P^{\perp}}
\newcommand{\Uex}{U_{\mathrm{ex}}}
\newcommand{\Vex}{V_{\mathrm{ex}}}
\newcommand{\Dex}{D_{\mathrm{ex}}}
\newcommand{\Dpre}{D_{\mathrm{pre}}}
\newcommand{\Djoint}{D_{\mathrm{joint}}}
\newcommand{\code}[1]{\texttt{\small #1}}
\newcommand{\file}[1]{\texttt{\small #1}}

\title{\textbf{RSVD-CBE as the basis update of the BUG integrator}\\[2mm]
\large Session record, method description, and a worked example}
\author{BUG-Julia, branch \code{feature/cbe-bug}}
\date{28 July 2026}

\begin{document}
\maketitle

\begin{abstract}
\noindent
The BUG (Basis-Update and Galerkin) integrator updates its bases by solving two auxiliary
ODEs per bond --- the $K$-step and the $L$-step --- and then takes one Galerkin step in the
augmented basis. This document describes replacing those two ODEs with \emph{controlled bond
expansion} (CBE) driven by a \emph{randomised SVD} (RSVD): instead of evolving a basis, one
applies $\Hb$ once, sketches the result against a random probe, and keeps the directions
carrying the largest weight. Part~\ref{part:session} records what was done, measured, and
\emph{retracted} in the session of 27--28 July 2026, including a wrong claim that survived
several hours and how it was caught. Part~\ref{part:method} explains the method: the RSVD
step from first principles, exactly which operations replace the $K$/$L$ ODEs, the knobs and
what each controls, and how the $S$-step is taken. Part~\ref{part:example} walks one bond
update numerically, printing every intermediate; all numbers there are output of
\file{examples/cbe\_walkthrough.jl} rather than illustration.
\end{abstract}

\tableofcontents

% =====================================================================================
\part{Session record}
\label{part:session}
% =====================================================================================

\section{Starting point}

The session opened with a handoff document reporting that RSVD-CBE worked in both integrator
paths after two defects in the selection had been fixed, and listing four remaining tasks. The
first act was to rename that handoff to \file{docs/current\_work.md}, so that one file states
what is true now rather than accreting a dated copy per session.

\section{What was decided and measured}

\subsection{The expansion budget: default \texorpdfstring{$\text{growth}=2.0$}{growth = 2.0}}
\label{sec:budget}

The expansion budget per bond was
\begin{equation}
  \Dex = \max\!\big(\lceil 1.1\, d_{\max}\rceil - r,\; 1\big),
  \qquad
  d_{\max} = \max(\chi_{i-1},\, r,\, \chi_{i+1}),
  \label{eq:budget}
\end{equation}
with $r$ the current bond rank; the constant $1.1$ is taken from the reference DMRG
implementation. Measurement showed this to be the accuracy limiter: on XX, $L=10$, the
quantity \code{err\_fnl} --- the fraction of ranked candidate weight discarded for want of
budget --- sat at $0.53$, and raising the budget improved the error from
$2.3\times10^{-4}$ to $1.2\times10^{-6}$.

The default is now $\text{growth} = 2.0$. The reasoning is not that $2.0$ is special but that
$1.1$ is a \emph{DMRG} schedule: there every sweep re-converges the same state, so a ten
percent ratchet compounds over sweeps and costs only iterations. A time integrator gets one
shot per step --- a direction the step needs and does not admit is not deferred, it is gone,
and the state carries that error forward. At $r=13$, Eq.~\eqref{eq:budget} proposes two
directions.

A ratio was chosen over an absolute floor (the measured saturation point was $\Dex=8$)
because $8$ is where \emph{that} chain saturates, whereas a ratio scales with the state. Wide
bonds remain bounded by \code{room} (the local space, a hard limit) and by \code{maxdim} at the
truncation, which is where a rank ceiling belongs.

\paragraph{The budget is problem-dependent.} This is recorded as a standing rule rather than a
caveat: the value was tuned on one chain (XX, $L=10$, domain wall, $t=0.5$), and a growth
schedule tracks how fast entanglement grows, how many channels $\Hb$ has, and how close
\code{maxdim} is to binding. Table~\ref{tab:decide} is the decision procedure, and it is cheap
to execute --- the $\Dex$ curve costs seconds per point at small $L$.

\begin{table}[t]
\centering
\begin{tabular}{@{}p{0.32\linewidth}p{0.34\linewidth}p{0.26\linewidth}@{}}
\toprule
\textbf{observation} & \textbf{meaning} & \textbf{action} \\
\midrule
\code{err\_fnl} $>0$ & the cap is the accuracy limiter & raise \code{growth} \\
\code{err\_fnl} $=0$, error still high & budget is not the limiter & look at $\delta t$,
\code{maxdim}, \code{comp\_ratio} \\
\code{err\_fnl} $=0$ and time matters & past saturation & lower \code{growth} until
\code{err\_fnl} leaves $0$ \\
\code{krylov\_dims} high, \code{mean\_aug} small & iterations dominate & raise \code{growth};
bigger can be cheaper \\
\code{krylov\_dims} flat, time rising & cost is in sketch/SVDs & lower \code{growth}, or pin
\code{dex} \\
\bottomrule
\end{tabular}
\caption{Choosing the budget for a new problem from recorded diagnostics.}
\label{tab:decide}
\end{table}

\subsection{Second order in \texorpdfstring{$\delta t$}{dt}, measured}

Whether the single centre Galerkin step costs an order in $\delta t$ was an open question in
the design notes, with first order named as the failure mode to watch for. It does not:
at $L=18$ with \code{maxdim}$=64$ (nothing truncates, \code{maxbond} peaks at 28),

\begin{center}
\begin{tabular}{@{}lrrrr@{}}
\toprule
$\delta t$ & CBE-BUG & ratio & 2-site TDVP & ratio \\
\midrule
$0.100$ & $1.0235\times10^{-3}$ & --- & $2.0640\times10^{-4}$ & --- \\
$0.050$ & $2.5648\times10^{-4}$ & $3.99$ & $5.1784\times10^{-5}$ & $3.99$ \\
$0.025$ & $6.4211\times10^{-5}$ & $3.99$ & $1.2969\times10^{-5}$ & $3.99$ \\
\bottomrule
\end{tabular}
\end{center}

\noindent
and $\mathrm{err}/\delta t^2$ is constant to three digits for CBE-BUG ($0.1023$, $0.1026$,
$0.1027$). CBE-BUG does carry a $\sim 5\times$ larger error \emph{constant} than 2-site TDVP at
equal $\delta t$ and equal untruncated rank; that is a statement about the subspace, not the
order, and the rank/memory comparison has to be read against it.

\subsection{A calibration that could not calibrate}

The campaign's $\delta t$ calibration ran only at \code{maxdim}$=64$, where \code{maxbond}
peaks at $28$: nothing truncates, so those numbers carry no rank information and cannot be used
to choose the main run's $\delta t$. A run at $\delta t = 0.05$ would have compared two
$\delta t$-limited runs and reported time-integration constants as if they were rank effects.
The calibration now takes \code{maxdim} and $t_{\max}$, so the choice is made where the main run
reports.

\subsection{The memory claim, plumbed}

\code{peak\_elements} --- the working set, \emph{state $+$ transients alive}, measured inside
the step at a definition both integrators share --- was already recorded by both integrators
but nothing reported it. The plots used \code{state\_elems} $+$ \code{theta\_elems} instead,
which flatters CBE-BUG twice: its \code{theta\_elems} is $0$ by construction, and its basis
sweep is transiently \emph{wider} than the state that survives truncation, so the state alone
understates its peak. The campaign now emits \code{peak\_elems} and the plots read it.

\section{The retraction}
\label{sec:retraction}

\textbf{A claim reported during this session was wrong: that a generous expansion budget makes
the integrator \emph{faster}.} It was measured as \code{gate} $24.4\,\mathrm{s} \to
3.9\,\mathrm{s}$ and \code{mpo} $19.7\,\mathrm{s}\to 3.6\,\mathrm{s}$, and it was an artefact
of measurement order: Julia compiles a call chain on first use, \code{@elapsed} charges that
compilation to the first timed row, and the starved setting ran first in every table.

What made it unarguable was a set of rows that are the \emph{same computation}. At
$\text{growth}=1.1$ the Lubich path stalls at \code{maxbond} $4$, so caps of $8$, $16$ and $32$
never bind:

\begin{center}
\begin{tabular}{@{}lrrrr@{}}
\toprule
run & error & \code{maxbond} & iterations/step & wall \\
\midrule
\code{chi<=8}  & $5.710\times10^{-3}$ & 4 & 4.9 & $160.6\,\mathrm{s}$ \\
\code{chi<=16} & $5.710\times10^{-3}$ & 4 & 4.9 & $4.6\,\mathrm{s}$ \\
\code{chi<=32} & $5.710\times10^{-3}$ & 4 & 4.9 & $4.6\,\mathrm{s}$ \\
\bottomrule
\end{tabular}
\end{center}

\noindent
Identical to four digits in every column except wall time, which differs by $35\times$. It also
reconciles the original numbers arithmetically: in the budget benchmark the validated
integrator ran first and absorbed most of the shared compilation, leaving the first CBE row
$\sim20\,\mathrm{s}$ of residue, so $24.4 \approx 3.9 + 20$ and $19.7 \approx 3.6 + 16$.

The hypothesis offered for the speedup --- that a starved basis makes the $0$-site exponential
harder to converge --- is also dead, and independently of any timing, because the counters are
collected inside the runs where contention cannot reach them: on the bond-update path the
starved and generous settings use $347.9$ versus $343.6$ iterations per step (identical), with
the generous run doing \emph{more} total work; on the Lubich path the starved setting uses
$6\times$ fewer iterations on $4\times$ smaller objects.

\paragraph{Fixes.} A \code{warmup} routine now runs one throwaway step per path and per
Hamiltonian before any timing; a \code{rev} mode reverses row order so an order-dependent
artefact cannot pass unnoticed; and two operational rules are recorded --- verify the machine is
idle before timing (killing a wrapper script does \emph{not} kill the Julia child it launched,
which cost two tables in one session), and never trust a single pass when rows that are the same
computation can be used as a consistency check.

\paragraph{The corrected cost picture}, warmed up and run in both row orders on a verified-idle
machine (XX, $L=10$, $\delta t = 0.02$, $t=0.5$, \code{maxdim}$=64$):

\begin{center}
\begin{tabular}{@{}lrrrrr@{}}
\toprule
setting & error & \code{maxbond} & forward & reversed & iters/step \\
\midrule
\code{cbe\_bond\_update} $1.1$ & $2.347\times10^{-4}$ & 14 & $10.5\,$s & $9.9\,$s & 347.9 \\
\code{cbe\_bond\_update} $2.0$ & $1.165\times10^{-6}$ & 13 & $9.5\,$s & $9.7\,$s & 343.6 \\
\code{cbe\_bond\_update} \code{maxiter=8} & $1.165\times10^{-6}$ & 13 & $6.2\,$s & $6.4\,$s & 105.8 \\
\code{cbe\_lubich} $1.1$ & $6.972\times10^{-6}$ & 4 & $3.5\,$s & $3.3\,$s & 4.8 \\
\code{cbe\_lubich} $2.0$ & $8.445\times10^{-6}$ & 14 & $6.8\,$s & $5.6\,$s & 27.8 \\
\code{cbe\_lubich} \code{maxiter=8} & $8.445\times10^{-6}$ & 14 & $6.1\,$s & $6.0\,$s & 7.6 \\
\bottomrule
\end{tabular}
\end{center}

\begin{itemize}[leftmargin=1.2em]
\item On \code{cbe\_bond\_update} the budget is \textbf{free} --- $1.0\pm0.1$ in either order ---
for a $200\times$ accuracy gain.
\item On \code{cbe\_lubich} the budget costs $\sim1.9\times$, the direction a cost model
predicts. The starved run is cheap only because it is stalled and not doing the physics: at
$L=12$, $t=1.5$ it is $22\times$ less accurate.
\item \code{maxiter}$=30$ is the wasteful default, not the budget --- capping at $8$ cuts
\code{cbe\_bond\_update} by $\sim36\%$ with the error unchanged to all printed digits, while
saving nothing on \code{cbe\_lubich}. The asymmetry is structural: one exponential per
\emph{bond} versus one per \emph{step}.
\end{itemize}

\section{Naming, and one further correction}

The two paths are now named for what they are, and every occurrence was renamed across sources,
tests, benchmarks and documents:

\begin{center}
\begin{tabular}{@{}ll@{}}
\toprule
old & new \\
\midrule
\code{cbe\_gate\_bond\_update} & \code{cbe\_bond\_update} \\
\code{cbe\_gate\_sweep!} & \code{cbe\_bond\_update\_sweep!} \\
\code{cbe\_gate\_bug!} & \code{cbe\_bond\_update\_bug!} \\
\code{cbe\_expand\_gate} & \code{cbe\_expand\_bond} \\
\code{sketch\_gate\_left/right} & \code{sketch\_bond\_left/right} \\
\code{cbe\_bug\_bond\_update} & \code{cbe\_lubich\_sweep} \\
\code{cbe\_bug!} & \code{cbe\_lubich\_bug!} \\
\bottomrule
\end{tabular}
\end{center}

\noindent
\textbf{Neither path applies a global $\Hb$ to the state}, and earlier wording in the documents
that said so was wrong. Every application is an \emph{effective local} operator: a two-site
effective $\Hb$ for the basis sketch, assembled from the left channel environments at link $i$,
the right channels at link $i+2$ and the local terms; and a zero-site effective $\Hb$ for the
Galerkin step, the environments contracted with the augmented frames. The only defensible sense
of ``global'' is \emph{which terms the effective operator contains}: all of $\Hb$ for
\code{cbe\_lubich\_sweep}, versus only the local bond term for \code{cbe\_bond\_update}, where
the rest of the chain reaches the bond solely through the canonical gauge. The environment
module says as much in its own header: \emph{not a symmetric MPO} --- it is an MPO contraction
with the virtual leg \emph{named} rather than fused.

\section{Verification}

Both suites pass with the new default and after the rename: $1019/1019$ in the CBE module suite
and $722/722$ in the validated integrator's suite, the latter untouched by this work.

% =====================================================================================
\part{The method}
\label{part:method}
% =====================================================================================

\section{Setting and notation}

Work at one bond $(i, i+1)$ of an MPS. Cutting there gives the two-site block
\begin{equation}
  \Theb = U_0\, S_0\, V_0,
\end{equation}
where $U_0$ carries legs $(\chi_{i-1}, d, r)$ and is a left isometry
($U_0^\dagger U_0 = \mathbb{1}$), $V_0$ carries $(r, d, \chi_{i+1})$ and is a right isometry,
$S_0$ is the $r\times r$ core, $d$ is the physical dimension and $r$ the bond rank. The
\emph{frames} $U_0, V_0$ define the variational manifold the step may move in; the whole
difficulty of a rank-adaptive integrator is choosing \emph{new columns} for them.

The two local spaces bound how far each side can grow:
\begin{equation}
  \mathrm{room}_l = d\,\chi_{i-1} - r,
  \qquad
  \mathrm{room}_r = d\,\chi_{i+1} - r .
  \label{eq:room}
\end{equation}
These are hard limits: a frame cannot have more columns than its local space has dimensions.

\section{What the $K$- and $L$-steps do, and why they can be replaced}
\label{sec:kl}

\subsection{The original $K$/$L$ steps}

In BUG the basis at a bond is updated by evolution. The $K$-step forms $K = U_0 S_0$ and
evolves it with the \emph{opposite} frame held fixed. In this codebase --- the
\emph{discarded} variant --- the complement projector is applied to the \emph{generator}, before
the exponential (\file{src/BondUpdateBUG/kls\_step.jl}):
\begin{equation}
  \dot K(t) = \Pperp_{U_0}\Big[\big(\Hb\,(K(t)\,V_0)\big)V_0^\dagger\Big],
  \qquad
  \Pperp_{U_0} = \mathbb{1} - U_0U_0^\dagger,
  \qquad K(0) = U_0 S_0,
  \label{eq:kstep}
\end{equation}
integrated over the step, after which one augments
$\hat U = \mathrm{orth}\big([\,K(\tau) \mid U_0\,]\big)$; the $L$-step is the mirror on the
right. Note the inner bracket: $\Hb$ acts on the two-site block $K V_0$ and the result is closed
against $V_0^\dagger$, so this is already an \emph{effective} operator, not a global one.
Applying $\Pperp$ to the generator makes it \textbf{non-Hermitian}, which is why the $K$/$L$
substeps take the Arnoldi path while only the $S$-step stays Hermitian.

Two observations drive everything that follows.

\begin{enumerate}[leftmargin=1.4em]
\item \textbf{Only the span survives.} $K(\tau)$ itself is discarded; what enters the Galerkin
step is the \emph{subspace} $\mathrm{range}([K(\tau) \mid U_0])$. Any procedure returning that
subspace is as good as solving Eq.~\eqref{eq:kstep}.
\item \textbf{To first order the span is an $\Hb$-image.} Expanding
Eq.~\eqref{eq:kstep},
\begin{equation}
  K(\tau) = U_0 S_0 - \mathrm{i}\tau\, \Pperp_{\text{right}} \Hb\, (U_0 S_0) + O(\tau^2),
\end{equation}
so
\begin{equation}
  \boxed{\;\mathrm{range}\big([\,K(\tau) \mid U_0\,]\big)
       = \mathrm{range}\big([\,U_0 \mid \Pperp (\Hb\Theb)\,]\big) + O(\tau).\;}
  \label{eq:equiv}
\end{equation}
\end{enumerate}

The right-hand side of Eq.~\eqref{eq:equiv} contains \emph{no time evolution}: one application
of $\Hb$, with the current frame projected out. That is the entire idea. The $\tau^2$ and higher
powers are not lost --- they are recovered by the Krylov iteration inside the $S$-step
(Sec.~\ref{sec:sstep}), which builds a genuine Krylov space over the \emph{expanded} basis.

\subsection{What replaces them, operation by operation}

\begin{center}
\begin{tabular}{@{}p{0.30\linewidth}p{0.32\linewidth}p{0.30\linewidth}@{}}
\toprule
& \textbf{standard BUG} & \textbf{RSVD-CBE} \\
\midrule
new left directions & solve Eq.~\eqref{eq:kstep} (\code{expv} on a
\emph{non-Hermitian} operator, Arnoldi) & one sketched application of $\Hb$, then two small
SVDs \\
new right directions & the mirror $L$-ODE & the mirror sketch \\
empty charge sectors & random per-sector seed columns (\code{missing\_fill}) & reached
\emph{through} $\Hb$ by a sector-graded probe \\
ranking of candidates & none: every direction found is admitted & candidates ranked by weight
in $\Hb\Theb$, top $\Dex$ kept \\
exponentials per bond & two (non-Hermitian) & \textbf{zero} \\
\bottomrule
\end{tabular}
\end{center}

\noindent
The last two rows are the substantive gains. Admitting every direction found means the rank
grows by whatever the ODE happened to produce; ranking means a small $\Dex$ can be
\emph{targeted}. And reaching an empty-but-reachable charge sector through $\Hb$ returns the
direction the dynamics actually wants, rather than a random column the $S$-step's SVD then
discards.

\paragraph{The $K$/$L$ machinery already needed an $\Hb\Theb$ patch.} This is the strongest
argument for the replacement, and it predates it. Because the $K$- and $L$-half-steps hold the
opposite frame fixed, on a product state the off-diagonal part of the generator in
Eq.~\eqref{eq:kstep} projects to zero and \emph{the rank cannot grow at all} --- the state
freezes, fatally so without $U(1)$ sectors where a random fill cannot help either. The validated
kernel therefore carries a ``criterion-2'' enrichment: it forms the residual of the \emph{full}
two-site action $\Hb\Theb$ and enriches the frames with the part of it orthogonal to what $K$/$L$
found. So the existing integrator already applies $\Hb$ once per bond and already grows its basis
from that residual --- it just does so as a corrective patch bolted onto two ODE solves.
RSVD-CBE keeps that mechanism, drops the ODEs and the random fill, and adds a weight ranking:
one $\Hb$ application doing the whole job it was previously doing only as a rescue.

\section{The RSVD step, from first principles}
\label{sec:rsvd}

\subsection{The problem it solves}

We need an orthonormal basis for the part of the left image of $\Hb\Theb$ that lies outside
$U_0$:
\begin{equation}
  \mathcal{C} \;=\; \mathrm{range}\big(\Pperp\,(\Hb\Theb)_{\text{left}}\big),
  \qquad \Pperp = \mathbb{1} - U_0 U_0^\dagger,
  \label{eq:target}
\end{equation}
and we want the $\Dex$ directions in $\mathcal{C}$ of largest weight. The direct route --- form
$\Hb\Theb$, matricise, take its SVD --- is exactly what this method exists to avoid:
$\Hb\Theb$ is the rank-4 two-site object that 2-site TDVP allocates and BUG does not.

\subsection{Randomised range finding in three sentences}

Let $A$ be a matrix whose column space we want but whose entries we would rather not form. Take
a thin random matrix $\Om$ with a few more columns than the rank we seek, and compute
$Y = A\,\Om$. Because a random vector has a component along every direction almost surely, the
columns of $Y$ span the dominant part of $\mathrm{range}(A)$; orthonormalising $Y$ therefore
gives a basis for it, and the ``few more columns'' (the \emph{oversampling}) is what makes the
statement reliable rather than merely likely.

That is the whole of RSVD as used here, with three specialisations.

\paragraph{(i) $A$ is never formed.} $A\Om$ is computed by folding $\Om$ into the operator
before contracting. In the bond-update path, $\Om$ is absorbed into the gate's bra leg first, so
the largest intermediate is $O(d^2 g r)$ and the rank-4 object never exists:
\begin{equation}
  Y = (\Hb\Theb)\,\Om^\dagger,
  \qquad
  \Om \;\text{in the layout of}\; V_0,\;\text{with } g \;\text{columns instead of } r .
\end{equation}
Setting $\Om = V_0$ recovers the exact left factor of $\Hb\Theb$, which is how the sketch is
tested against a formed reference.

\paragraph{(ii) The probe is drawn per charge sector.} With a $U(1)$ symmetry the bond leg
carries charge sectors, and a probe drawn without regard to them can miss a sector entirely ---
including a sector that is currently empty but reachable by $\Hb$. The probe's columns are
therefore allocated across sectors in proportion to sector dimension. This is what makes the
random per-sector seeding of standard BUG unnecessary.

\paragraph{(iii) The probe is half complement, half isometry.} The probe is built in two halves:
random columns with $V_0$ projected \emph{out} (the complement half), and random columns with
$V_0$ projected \emph{in} (the isometry half), mixed by \code{comp\_ratio} (default $0.5$). It
is worth being precise about what this does \emph{not} do: it does not control what is admitted.
The admitted block is always outside the frame, because the preselection projects the frame out
regardless. \code{comp\_ratio} $=1$ therefore does not mean ``keep only the discarded space'';
it means the probe is blind to the one-site component, which the reference implementation warns
matters, and it caps the probe's width at the opposite frame's complement dimension --- which
near a chain edge is $0$ or $1$.

\subsection{Budget, oversample, probe, candidates: the four steps in code}
\label{sec:fourfunnel}

The four quantities are produced in a strict order, each from the one before, and the
distinction between them is where most of the confusion about this method lives. All line
references are \file{src/RSVDCBEBondUpdate/cbe.jl}.

\paragraph{Step 1 --- the budget $\Dex$: \emph{how many new directions this bond may admit}.}
\begin{lstlisting}[style=jl]
dmax         = max(leg_dim(U0,1), r, leg_dim(V0,3))          # neighbourhood max
budget       = dex <= 0 ? max(ceil(Int, growth*dmax) - r, 1) : dex
dex_l, dex_r = min(budget, room_l), min(budget, room_r)      # room = d*chi - r
\end{lstlisting}
Read it in four moves. \code{dex} (default $0$) is an \emph{absolute override}: any positive
value bypasses the schedule completely, which is what benchmarks use so the budget is a fixed
axis. Otherwise \code{growth} (default $2.0$) sets it relative to $d_{\max}$ --- the
\textbf{neighbourhood} max over the left, centre and right bond dimensions, not $r$ alone, so a
wide bond can pull narrow neighbours up. The \code{max(..., 1)} floor guarantees every bond may
propose at least one direction. Finally each side is clipped by its own
$\mathrm{room} = d\chi - r$, a hard ceiling: no frame can exceed the local space it lives in. If
both sides clip to zero the routine returns immediately --- the frames are already complete, and
that is success, not failure.

\paragraph{Step 2 --- the oversample $\Dpre$: \emph{how wide the random probe is}.}
\begin{lstlisting}[style=jl]
npre = preselect_only ? max(budget, 1) :
       dover === nothing ? ceil(Int, 1.2 * budget) : budget + dover
\end{lstlisting}
Three modes. Default (\code{dover = nothing}) is \textbf{relative}: $\Dpre = \lceil 1.2\,
\text{budget}\rceil$, the reference's $20\%$ margin. Passing \code{dover = n} makes it
\textbf{additive}, $\Dpre = \text{budget} + n$, which gives small budgets proportionally more
help. \code{preselect\_only} drops the margin entirely, since in that mode there is no ranking
stage for the extra columns to feed.

\emph{Why oversample at all.} A probe of exactly $k$ columns does span a $k$-dimensional slice of
$\mathrm{range}(A)$, but the chance that this particular random slice is well aligned with the
true top-$k$ subspace is not close to $1$ --- the projection can be arbitrarily ill-conditioned.
A handful of extra columns is what turns ``probably captures the range'' into ``misses it with
exponentially small probability''. It is a reliability margin, not an accuracy parameter.

\emph{One subtlety worth flagging}, because it is invisible in the code at a glance: \code{npre}
is computed from the \textbf{unclipped} \code{budget}, not from \code{dex\_l}/\code{dex\_r}. A
bond with almost no room still gets a full-width probe.

\paragraph{Step 3 --- the probe $\Om$ itself (\code{sector\_graded\_sketch}).}
Given $\Dpre$, build a tensor in the \emph{opposite} frame's layout:
\begin{enumerate}[leftmargin=1.4em,itemsep=1pt]
\item \code{fusion\_basis} gives the full local space; per charge sector $q$ it has dimension
$d_{\text{full}}(q)$. The frame already spans $d_T(q)$ of it, leaving a complement
$d_C(q) = d_{\text{full}}(q) - d_T(q)$.
\item Split the width: $n_c = \mathrm{round}(\Dpre \cdot \code{comp\_ratio})$ complement columns,
$n_t = \Dpre - n_c$ isometry columns.
\item \code{\_alloc\_columns} spreads each half across sectors \emph{in proportion to sector
dimension}, capped at each sector's own dimension, using a \textbf{per-sector \code{ceil}}.
\item Any shortfall --- a half whose subspace is too small to supply its quota --- is
\textbf{handed to the other half}, and back again if that one is short too.
\item The complement half is random-then-projected-\emph{out} of the frame, the isometry half
random-then-projected-\emph{in}; the two are direct-summed.
\end{enumerate}

Two consequences that surprise people, both load-bearing:

\begin{itemize}[leftmargin=1.2em]
\item \textbf{$\Dpre$ is a target, not the realised width.} The \code{ceil} in step 3 is applied
\emph{per sector}, so the sum over sectors can exceed $\Dpre$. In the worked example
$\Dpre = 9$ yet the probe comes out $10$ columns wide, split
$[-2 \Rightarrow 2,\; 0 \Rightarrow 4,\; 2 \Rightarrow 3,\; 4 \Rightarrow 1]$. This is
deliberate: the \code{ceil} guarantees every sector with room receives at least one column,
which is precisely what matters when the sector you need is the currently \emph{empty} one. A
\code{floor} would round the small sectors to zero and the probe would never reach them.
\item \textbf{The redistribution in step 4 is not tidiness.} Near a chain edge the opposite
frame's complement is $0$- or $1$-dimensional. Without redistribution a complement-weighted
split collapses the probe to one column and the bond declines to expand \emph{with room to
spare} --- measured, leaving right residuals of $3.1\times10^{-4}$ and $0.19$ untouched, and one
blocked narrow bond then caps every bond inward.
\end{itemize}

\paragraph{Step 4 --- the preselection (\code{\_preselect\_left}).}
\begin{lstlisting}[style=jl]
R = perp_component(U0, Y)                             # P_perp Y, on legs (link_l, site_l)
n = norm(R)
n <= 100*eps * max(norm(Y), 1.0) && return nothing    # nothing outside the frame
res = svd((1/n) * R, (1,2); cutoff = stol_pre)        # NORMALISED first
Q   = res.U
R2  = perp_component(U0, Q)                           # RepairOrtho
Q   = svd(R2, (1,2); cutoff = 1e-13).U
return (Q, trunc_weight(res))                         # the tail becomes err_pre
\end{lstlisting}
Project the frame out; bail if what remains is noise \emph{relative to} $\|Y\|$; normalise;
SVD at \code{stol\_pre} ($10^{-10}$); then orthonormalise a \emph{second} time. The
normalisation and the second pass are the two steps that look redundant and are not --- both are
explained in \S\ref{sec:twostage} and in Stage~4 of the worked example.

The output width is the \textbf{number of candidates}, and it is decided by the SVD cutoff alone
--- not by $\Dex$, not by $\Dpre$. This is the point most easily got wrong, so state it flatly:

\begin{center}
\small
\begin{tabular}{@{}p{0.30\linewidth}p{0.62\linewidth}@{}}
\toprule
\textbf{Does the budget affect preselection?} & \textbf{Only through the probe's width.} \\
\midrule
Directly & No. Preselection is governed by \code{stol\_pre} and by what $\Hb$ actually
produces outside the frame. \\
Indirectly & Yes, and decisively: the budget sets $\Dpre$, and the probe can never yield more
candidates than it has columns. A budget too small starves the probe and the directions are
never \emph{generated}. \\
Where it re-enters & At the \emph{final} selection, where \code{dex\_l}/\code{dex\_r} cap how
many candidates are admitted, and \code{err\_fnl} reports what the cap cost. \\
\bottomrule
\end{tabular}
\end{center}

\noindent
That split is the whole design. A starved budget fails \emph{silently at step 3} --- there is no
diagnostic for a direction that was never probed. A generous budget lets the preselection decide
on its own merits and leaves the cap non-binding, which is the \code{err\_fnl}$\,=0$ healthy
case. It is why the $\text{growth}=1.1$ schedule showed \code{err\_fnl}$\,=0.62$: not a
selection failure, a \emph{budget} failure.

\paragraph{The funnel, on the worked example.} Every number below is from the transcript in
Part~\ref{part:example}, at bond $(4,5)$ of an $L=8$ chain with $r=7$, $\chi_l=\chi_r=6$, $d=2$.

\begin{center}
\small
\begin{tabular}{@{}llll@{}}
\toprule
quantity & symbol & value & set by \\
\midrule
neighbourhood max        & $d_{\max}$ & $7$    & $\max(6,7,6)$ \\
budget                   &            & $7$    & $\lceil 2.0\cdot 7\rceil - 7$ \\
room, each side          &            & $5,\,5$ & $d\chi - r$ \\
per-side budget          & $\Dex$     & $5$    & $\min(7, 5)$ \\
oversampled target       & $\Dpre$    & $9$    & $\lceil 1.2\cdot 7\rceil$ \\
\emph{realised} probe width &         & $10$   & per-sector \code{ceil} \\
candidates surviving preselection &   & $4$    & SVD at \code{stol\_pre} \\
admitted                 &            & $4$    & $\min(4, 5)$ --- cap does not bind \\
rank after expansion     &            & $7 \to 11$ & direct sum \\
rank after the $S$-step  &            & $10$   & truncation \\
\bottomrule
\end{tabular}
\end{center}

\noindent
Read down that table and the method is one sentence: a budget of $7$ becomes a probe of $10$
columns, which finds $4$ genuinely new directions, all $4$ of which fit under the cap of $5$ --- and
under the old schedule the budget would have been $1$, the probe $2$ columns, and three of those
four directions would never have existed to be ranked.

\subsection{The two-stage selection}
\label{sec:twostage}

\begin{description}[leftmargin=1.4em,style=nextline]
\item[Stage 1 --- preselection (cheap, generous).]
Sketch, project the frame out, orthonormalise:
\begin{equation}
  Y = (\Hb\Theb)\Om^\dagger,
  \qquad
  R = \Pperp Y,
  \qquad
  Q_L = \mathrm{orth}(R) \;\text{via SVD at tolerance \code{stol\_pre}} .
\end{equation}
$R$ is normalised before the SVD because the cutoff is relative to the largest singular value
and the new directions are $O(\tau)$-small in absolute terms. The orthonormalisation is then
\emph{repeated} --- the reference's \code{RepairOrtho} --- because a direction surviving the
cutoff can still carry a component along $U_0$ of order $\varepsilon/\text{cutoff}$, and the
direct sum $[U_0 \mid Q_L]$ is an isometry only if $Q_L \perp U_0$ exactly. This is not
bookkeeping: the same omission elsewhere in the codebase grew a state's norm by $128\times$
over five steps.

\item[Stage 2 --- final selection (the ranking).]
Form the \emph{small} matrix
\begin{equation}
  M = Q_L^\dagger\,(\Hb\Theb)\,Q_R
  \qquad (\Dpre \times \Dpre),
  \label{eq:M}
\end{equation}
take its SVD and keep $\Dex$ singular directions. Its singular values \emph{are} the weights
of the candidate pairs in the two-site action, so this is the ranking step, and it is what the
$K$/$L$ augmentation structurally cannot do.
\end{description}

\noindent
One subtlety that was a defect and is now the design: this selection is applied \emph{per side}.
The reference's CBE is directional --- it expands one side per call --- so its joint matrix $M$
only ever ranks the side being expanded. Reusing $M$ to rank both sides at once lets a saturated
side veto its partner, and it did: capping each side at $\min(\Dex^{(l)}, \Dex^{(r)})$ throttled
one side to the other's headroom, and a global ``$M$ is negligible'' verdict returned nothing on
\emph{either} side while a residual of $0.19$ went unaddressed on one of them. Each side now
decides independently, falling back to its own preselected candidates when the joint SVD cannot
serve its whole budget --- ordered by the preselection SVD's singular values, though see the
caution in \S\ref{sec:clcr} on how that ordering interacts with charge sectors.

\subsection{Constructing the appended blocks \texorpdfstring{$C_L$}{CL} and \texorpdfstring{$C_R$}{CR}}
\label{sec:clcr}

$C_L$ and $C_R$ are the candidate directions after ranking and after trimming to budget. They are
built \emph{per side and independently} (\code{cbe\_expand}), and each side takes one of two
branches. Write $\Djoint = \min(\Dex^{(l)}, \Dex^{(r)}, \Dpre^{(l)}, \Dpre^{(r)})$ for the number
of directions the joint SVD can rank.

\begin{description}[leftmargin=1.4em,style=nextline]
\item[Ranked branch.] Take the SVD $M = U_M \Sigma_M V_M^\dagger$ of the small matrix
Eq.~\eqref{eq:M}, truncated at \code{stol\_fnl} to $\Djoint$ directions, and \emph{rotate the
preselected candidates into that ranked basis}:
\begin{equation}
  C_L \;=\; Q_L\,U_M,
  \qquad
  C_R \;=\; V_M^\dagger\,Q_R .
  \label{eq:clcr}
\end{equation}
No new directions are created here --- $C_L$ spans exactly the same space as $Q_L$; the rotation
only \emph{orders} that space so that trimming to $\Dex$ keeps the heaviest combinations.

\item[Fallback branch.] Otherwise the side keeps the first $\Dex$ columns of its own preselected
$Q$ (\code{\_trim\_total}), with no joint rotation at all.
\end{description}

\noindent
The branch condition is \code{dex\_l <= joint}, and since $\Djoint \le \Dex^{(l)}$ by
construction, that fires \emph{only} when $\Djoint = \Dex^{(l)}$ --- i.e.\ only when the joint SVD
can serve that side's entire budget. This is the ``each side decides independently'' design of the
previous subsection expressed in code: a side whose budget the joint ranking cannot fill does not
get a partially-ranked block, it falls back to its own candidates.

Two properties matter and both are inherited rather than re-established. First,
\textbf{orthonormality}: $Q_L$ is orthonormal and $\perp U_0$ (the \code{RepairOrtho} pass), and
$U_M$ is an isometry, so $C_L$ is orthonormal and $\perp U_0$ --- which is precisely what makes
$[\,U_0 \mid C_L\,]$ an isometry. The right side is the mirror image, with $V_M^\dagger$ carrying
orthonormal rows. Second, \textbf{tagging}: the new block's bond leg is retagged to $U_0$'s bond
tag before the direct sum, because the two preselections deliberately use distinct tags and the
tensor library forbids duplicate indices.

\begin{center}
\small
\begin{tabular}{@{}ll@{}}
\toprule
\textbf{Caution} & \textbf{Why} \\
\midrule
The fallback is not globally weight-ordered
  & \code{\_trim\_total} walks charge sectors in storage order and takes \\
  & a prefix of each, filling greedily. Within a sector the columns are \\
  & SVD-ordered, but the \emph{allocation across} sectors is by storage \\
  & order, not by weight. A heavy direction in a late sector can lose \\
  & to a light one in an early sector. Unquantified --- worth a test. \\
\bottomrule
\end{tabular}
\end{center}

\subsection{The expansion is lossless}

The new columns are appended by direct sum,
\begin{equation}
  \Uex = [\,U_0 \mid C_L\,],
  \qquad
  \Vex = [\,V_0 \mid C_R\,],
\end{equation}
and the core is re-expressed in the wider basis,
$\hat S_0 = \Uex^\dagger \Theb \Vex^\dagger$. Since $\Uex$ contains $U_0$ as a block,
\begin{equation}
  \Uex \hat S_0 \Vex = \Theb
  \qquad\text{exactly},
\end{equation}
so the expansion changes \emph{no amplitude} --- it only widens the space the step may use, and
$\hat S_0$ is identically zero on the new directions. Two consequences that repeatedly caused
confusion, and are worth stating plainly:

\begin{itemize}[leftmargin=1.2em]
\item \textbf{An expansion cannot raise the Schmidt rank.} A widened bond collapses back at the
next gauge move, and that is correct reporting rather than a lost expansion. Rank grows only
where the state is \emph{evolved} in the widened space.
\item \textbf{Therefore the evolution must happen before the truncation}, which it does: the
$S$-step makes $\hat S_1$ dense on the new directions, and only then is the bond compressed.
\end{itemize}

\section{The $S$-step}
\label{sec:sstep}

The $S$-step is unchanged from the validated integrator, and deliberately so. In the expanded
basis, define the projected generator
\begin{equation}
  \mathcal{A}(X) \;=\; \Uex^\dagger \,\Hb\, \big(\Uex X \Vex\big)\, \Vex^\dagger,
\end{equation}
and take one exponential,
\begin{equation}
  \hat S_1 = \exp(\tau \mathcal{A})\, \hat S_0,
  \qquad \tau = -\mathrm{i}\,\delta t ,
\end{equation}
by Lanczos, matrix-free, with $\mathcal{A}$ supplied as a closure. Then SVD $\hat S_1$, truncate
to \code{maxdim} and \code{trunc\_thresh}, and write the two cores back.

Three properties matter.

\begin{enumerate}[leftmargin=1.4em]
\item \textbf{Exact at full rank.} When the expanded frames span the whole local space,
$\mathcal{A}$ \emph{is} the local Hamiltonian and the step is the exact two-site propagator.
This is the first thing the tests check.
\item \textbf{$\tau = 0$ is the identity}, because $\hat S_0$ is lossless.
\item \textbf{Hermitian, hence Lanczos}, not Arnoldi --- unlike the $K$/$L$ ODEs it replaces,
whose generator $\Pperp\Hb$ is non-Hermitian. Removing them removes a non-symmetric Krylov
problem from the integrator.
\end{enumerate}

\paragraph{Where the two paths differ.} \code{cbe\_bond\_update} takes one such step \emph{per
bond} inside an odd/even Trotter sweep, with the bare gate as $\Hb$; the projection
$\Uex X \Vex$ builds a rank-4 block on every iteration, which is affordable at small $\chi$ and
is why this path is the diagnostic vehicle rather than the production one.
\code{cbe\_lubich\_sweep} takes \emph{one} step per time step, at the centre bond, with $\Hb$
dressed by channel environments; its $\mathcal{A}$ contracts a rank-2 core against rank-2 and
rank-3 environment blocks, so no rank-4 object is ever formed. This is the measured origin of the
\code{maxiter} asymmetry in Sec.~\ref{sec:retraction}: $\sim344$ iterations per step in the
first case, $\sim28$ in the second.

\section{Cost}

With $\Dpre \sim 1.2\,\text{growth}\cdot\chi$, per bond:

\begin{center}
\begin{tabular}{@{}llll@{}}
\toprule
stage & cost & with $\Dpre \propto \chi$ & note \\
\midrule
sketch $(\Hb\Theb)\Om^\dagger$ & $O(w\,d\,\chi^2 \Dpre)$ & $O(\chi^3)$ & $w$ = channel count \\
preselection SVD & $O(d\,\chi\, \Dpre^2)$ & $O(\chi^3)$ & thin \\
final SVD & $O(\Dpre^3)$ & $O(\chi^3)$ & small matrix \\
$S$-step, per iteration & $O(w\,\chi^2)$ & $O(\chi^2)$ & rank-2 only (Lubich path) \\
\bottomrule
\end{tabular}
\end{center}

\noindent
So the ratio schedule costs the same \emph{order} as the canonicalisation the sweep already
pays --- a larger constant, not a worse exponent. A fixed $\Dex$ would be $O(\chi^2 \Dex)$,
asymptotically cheaper, which is precisely the trade the accuracy measurements rejected.

Measured at small $\chi$ ($L=12$, caps $8/16/32$) the fitted exponent is far below cubic
($k \approx 0.05$ then $0.43$ per doubling) because $L$-dependent constants dominate there; the
asymptotic figure needs larger $\chi$ than a laptop should be asked for. No super-cubic
behaviour appears.

\paragraph{What can blow up, in order of how real it is.}
\begin{enumerate}[leftmargin=1.4em]
\item \textbf{Norm growth from non-orthogonal direct sums} --- the one that actually happened
($128\times$ over five steps, elsewhere in the codebase). Guarded by the repeated
orthonormalisation.
\item \textbf{A strict SVD selector} dropping exactly-zero directions, so re-deriving a frame
from a zero-padded tensor collapses it. This caused an early rank freeze; the fix was to
accumulate frames separately and assemble the state only after the step.
\item \textbf{What cannot blow up:} there is no inverse anywhere in this path --- no $S^{-1}$,
no overlap matrices to invert --- so ill-conditioned small singular values have nothing to
divide into. Rank is bounded by Eq.~\eqref{eq:room} and by \code{maxdim}, and the expansion is
lossless.
\end{enumerate}

\section{The knobs}

\begin{center}
\begin{tabular}{@{}llp{0.44\linewidth}@{}}
\toprule
knob & default & what it controls \\
\midrule
\code{growth} & $2.0$ & the budget schedule
$\Dex = \lceil \text{growth}\cdot d_{\max}\rceil - r$, neighbourhood-coupled through
$d_{\max}$. Problem-dependent; see Table~\ref{tab:decide}. \\
\code{dex} & $0$ & absolute per-side budget; any positive value \emph{overrides}
\code{growth} entirely. Use for measurement, where a fixed axis is wanted. \\
\code{dover} & \code{nothing} & oversampling of the probe; \code{nothing} gives
$\Dpre = \lceil 1.2 \Dex\rceil$. This is the RSVD reliability margin. \\
\code{comp\_ratio} & $0.5$ & complement/isometry split of the \emph{probe}. Does not control
what is admitted. Leave alone unless testing. \\
\code{stol\_pre} & $10^{-10}$ & relative cutoff of the preselection SVD. \\
\code{stol\_fnl} & $10^{-13}$ & relative cutoff of the final selection SVD. \\
\code{maxiter} & $30$ & Lanczos budget of the $S$-step. Measured wasteful on
\code{cbe\_bond\_update} ($\sim36\%$ of its time for no accuracy); irrelevant on
\code{cbe\_lubich}. \\
\code{maxdim}, \code{trunc\_thresh} & $200$, $10^{-12}$ & the truncation after the $S$-step ---
where a rank ceiling belongs. \\
\code{sulz\_cap} & \code{false} & optional global $2r_{\max}$ bound; belongs to the
rank-adaptive BUG, kept only for an A/B. \\
\code{preselect\_only} & \code{false} & skip the ranking and take preselected candidates
directly (the reference's \code{-p} mode). Diagnostic: isolates what the ranking is worth. \\
\code{expand} & \code{true} & \code{false} freezes the frames, giving a fixed-rank projector
step. \\
\bottomrule
\end{tabular}
\end{center}

% =====================================================================================
\part{A worked example}
\label{part:example}
% =====================================================================================

\noindent
Everything below is the output of \file{examples/cbe\_walkthrough.jl}: XX chain, $L=8$, $U(1)$
symmetric, a domain wall warmed by three steps of $\delta t = 0.05$ so the bond carries real
structure, then one bond update at bond $(4,5)$. Run it with

\begin{lstlisting}[language=bash]
julia --project=. examples/cbe_walkthrough.jl
\end{lstlisting}

\noindent
The bond is \textbf{chosen, not assumed}: the script takes the bond maximising
$\min(\mathrm{room}_l, \mathrm{room}_r)$. That matters, and the first draft of this example got
it wrong --- the centre cut of a short chain \emph{saturates} (at $L=6$ it caps at rank $8$ and a
warmed state reaches it), giving $\mathrm{room} = 0$ on both sides, so CBE correctly declined to
expand and every selection stage printed empty. A saturated bond is not a failure of the method
and must not be read as one.

The stages below are the stages of the method. Each is stated as: what is computed,
what came out, and what the number \emph{means} --- because most of these quantities are only
interesting relative to something else (a residual relative to the action, a discarded weight
relative to the kept spectrum), and a bare figure invites the wrong reading.

\section{Stage 0 --- the bond before anything happens}

Three warm-up steps give the state \code{bond\_dims = [2, 4, 6, 7, 6, 4, 2]}. Cutting at bond
$(4,5)$, where $r=7$ and both halves have five dimensions in hand, gives the frames, the core,
and the two quantities that bound everything afterwards --- $\mathrm{room}_l$ and
$\mathrm{room}_r$ of Eq.~\eqref{eq:room}, both $5$ here.

\begin{lstlisting}
chain          L = 8, XX, U(1), domain wall, 3 steps of dt = 0.05 (t = 0.15)
state          bond_dims = [2, 4, 6, 7, 6, 4, 2]
bond           (4, 5) -- chosen as the most expandable, rank r = 7
U0  (6 x 2 x 7)       left frame   (link_l, site_l, bond)
S0  (7 x 7)           core
V0  (7 x 2 x 6)       right frame  (bond, site_r, link_r)
Theta0 (6 x 2 x 2 x 6)  the two-site block, U0*S0*V0
bond sectors on U0 leg 3: [-4 => 1,  -2 => 3,  0 => 3]
left isometry defect  = 1.02e-15   (U0 is an isometry)
room_l = d*chi_l - r = 6*2 - 7 = 5
room_r = d*chi_r - r = 2*6 - 7 = 5
\end{lstlisting}

\noindent
The left isometry defect confirms $U_0^\dagger U_0 = \mathbb{1}$ to machine precision, which is
what makes the direct sum in Stage~6 an isometry rather than merely a wider matrix. The bond leg
carries several $U(1)$ charge sectors; the probe in Stage~2 has to respect that decomposition,
which is the whole reason for a \emph{sector-graded} sketch.

\section{Stage 1 --- the budget}

The two schedules are evaluated side by side on the same bond, which is the cleanest possible
statement of what the session's default change did.

\begin{lstlisting}
growth = 1.1 :  dmax = 7 -> budget = ceil(1.1*7) - 7 = 1    npre = 2
growth = 2.0 :  dmax = 7 -> budget = ceil(2.0*7) - 7 = 7    npre = 9

using growth = 2.0:  budget = 7, clipped per side to (dex_l, dex_r) = (5, 5)
probe width npre = ceil(1.2*budget) = 9      (the 1.2 is the RSVD oversampling)
\end{lstlisting}

\noindent
Read the two rows against each other: at the same $d_{\max}$, the reference's DMRG ratchet and
the new default differ by roughly an order of magnitude in the number of directions the bond is
permitted to propose. The probe width $\Dpre = \lceil 1.2 \Dex\rceil$ follows the budget, so the
schedule sets not only what is admitted but how wide a net is cast to find it --- which is why a
starved schedule cannot be compensated for downstream.

\section{Stage 2 --- the random probe (the R in RSVD)}

\begin{lstlisting}
Omega_R (10 x 2 x 6)      a stand-in RIGHT frame with 10 columns instead of r = 7
its columns are drawn PER CHARGE SECTOR: [-2 => 2,  0 => 4,  2 => 3,  4 => 1]
half the columns are random-then-projected-OUT of V0 (the complement half),
half are random-then-projected-IN  (the isometry half, comp_ratio = 0.5).
\end{lstlisting}

\noindent
$\Om$ is a \emph{stand-in} for the right frame: same leg structure, a different number of
columns. Its per-sector counts are the point --- a probe drawn without regard to charge would
leave whole sectors unsampled, and a sector currently empty but reachable by $\Hb$ is exactly
where the new physics lives.

Note that $\Om$ has $10$ columns where $\Dpre = 9$ was asked for. That is not an error: the
allocator takes $\lceil\,\cdot\,\rceil$ per sector, so that a sector entitled to a fractional
share still receives a whole column, and the total can exceed the request by less than the number
of sectors. Also note the probe's charge sectors ($-2,0,2,4$) differ from the bond's
($-4,-2,0$): the probe lives in the \emph{fused} local space of the right half, which is a
larger set of charges than the current bond carries --- and reaching those is the mechanism by
which a new sector can enter the state at all.

\section{Stage 3 --- one application of $\Hb$, sketched}

\begin{lstlisting}
Y   (6 x 2 x 10)      the left image of the gate's action, 10 columns wide
check: sketch with Omega = V0 equals contracting the formed H*Theta:  2.58e-22
H*Theta itself is rank 4 (6 x 2 x 2 x 6) and IS NEVER FORMED in the method;
it is built here only to check the sketch. Omega is folded into the gate first.
\end{lstlisting}

\noindent
Two things to note. First, the check: sketching with $\Om = V_0$ reproduces the result of forming
$\Hb\Theb$ and contracting it, to machine precision --- so the sketch is the exact object,
merely evaluated in an order that never materialises the rank-4 tensor. Second, that rank-4
tensor is built \emph{here and only here}, to perform that check; the method itself never
allocates it.

\section{Stage 4 --- preselection}

The candidate directions are the part of the sketched action lying outside the current frame.

\begin{lstlisting}
||Y|| = 1.017861e+00     ||P_perp Y|| = 7.945719e-06   ->  ratio 7.806e-06
Q_L (6 x 2 x 4)       4 orthonormal candidate directions, all orthogonal to U0
weight the preselection SVD discarded: err_pre = 1.696e-12
Q_L is orthogonal to U0 to 3.77e-17  (the RepairOrtho second pass earns this)
\end{lstlisting}

\noindent
\textbf{The most instructive number in the whole walkthrough is
$\|\Pperp Y\|/\|Y\| \approx 8\times10^{-6}$.} The part of the gate's action lying outside the
current frame is a \emph{millionth} of the action's norm --- and that tiny sliver is exactly the
new physics, the thing the basis update exists to capture. It is small for a structural reason:
the frame was built by the previous step, so it already spans almost everything $\Hb$ does to
the block; what is new is $O(\tau)$ and, here, $O(\tau^2)$-ish in norm.

Two consequences follow, and both are visible in the implementation.

\begin{itemize}[leftmargin=1.2em]
\item \textbf{The preselection must normalise before its SVD.} Telum's cutoff is relative to the
largest singular value of the matrix it is given. Handing it $\Pperp Y$ unnormalised, with a
tolerance meaningful on the scale of $\|Y\|$, would discard the entire complement as numerical
noise. \code{\_preselect\_left} therefore divides by $\|\Pperp Y\|$ first, so the tolerance is
relative to the new directions' \emph{own} scale.
\item \textbf{Absolute thresholds anywhere in this path are a bug.} A cutoff of, say, $10^{-8}$
applied to the action would delete every candidate at this bond.
\end{itemize}

\noindent
The last line of the block is the \code{RepairOrtho} pass earning its place: the candidates are
orthogonal to $U_0$ to $4\times10^{-17}$, without which the direct sum would not be an isometry
and the norm would drift.

\section{Stage 5 --- the ranking}

\begin{lstlisting}
UMU (4 x 4)          = Q_L^dag (H Theta) Q_R   -- a SMALL matrix, 4 x 4
singular values KEPT  (4): 3.039e-11  2.291e-12  4.951e-13  2.790e-15
singular values CUT   (0): -- none --
These singular values ARE the weights: this is the ranking the K/L augmentation
cannot do, since it admits every direction it finds regardless of weight.
\end{lstlisting}

\noindent
This is the step the $K$/$L$ augmentation cannot perform. The singular values of
$M = Q_L^\dagger (\Hb\Theb) Q_R$ \emph{are} the weights of candidate pairs in the two-site
action, so keeping the top $\Dex$ is a principled choice rather than an arbitrary cut. Note the
size of $M$: $4\times4$, and its SVD is the cheapest part of the whole procedure.

Their absolute magnitudes ($10^{-11}$ down to $10^{-15}$) are again the $O(\tau)$ smallness of
Stage~4, not a sign that the candidates are worthless; what matters is their \emph{ordering} and
their spread --- four orders of magnitude between the first and the last, so the ranking is
carrying real information about which directions to prefer.

Nothing is cut here, and it is worth being precise about why: there are $4$ candidates and the
budget permits $5$, so the cap does not bind at this bond. This is the \code{err\_fnl}$\,=0$
condition of Table~\ref{tab:decide} --- the healthy case. Under the old $\text{growth} = 1.1$
schedule the budget would have been $1$ (Stage~1), three of these four directions would have been
discarded, and \code{err\_fnl} would have reported the loss.

\paragraph{A caveat this transcript makes easy to misread.} The ranking is \emph{computed and
printed} above, but at this bond it is \emph{not} what builds the appended block. Following
\S\ref{sec:clcr}: $\Dex^{(l)} = 5$ while $\Djoint = \min(5,5,4,4) = 4$, so the guard
$\Dex^{(l)} \le \Djoint$ is $5 \le 4$, which is false --- \emph{both} sides take the fallback
branch, and since each $Q$ has $4$ columns against a budget of $5$, the trim returns them
untouched. So here
\begin{equation}
  C_L = Q_L, \qquad C_R = Q_R \qquad\text{(no rotation by } U_M, V_M^\dagger\text{)},
\end{equation}
and the joint SVD's only effect is to report \code{err\_fnl}$\,=0$. The column count cannot
distinguish the two branches --- both would yield $4$ --- so this is visible from the branch
condition alone, not from the output. The ranked branch fires when a bond is
\emph{candidate-rich relative to its budget}, which is the regime a tighter \code{growth} or a
wider bond puts you in; it is exactly the regime the old $1.1$ schedule lived in, and where the
$62\%$ discarded weight of \S\ref{sec:budget} was being thrown away.

\section{Stage 6 --- the expansion, and its losslessness}

\begin{lstlisting}
U_ex (6 x 2 x 11)       r = 7 -> 11   (+4 new left directions)
V_ex (11 x 2 x 6)       r = 7 -> 11   (+4 new right directions)
err_pre = 5.322e-12   err_fnl = 0.000e+00   (err_fnl = 0 means nothing ranked out)
LOSSLESS: ||U_ex S_hat0 V_ex - Theta0|| = 1.58e-17     ||S_hat0|| - ||Theta0|| = 0.00e+00
So the expansion changes NO amplitude -- it only widens the space the step may use.
This is also why an expansion alone cannot raise the Schmidt rank.

left residual ||HT - U U^dag HT|| / ||HT||:  before 7.785e-06  ->  after 2.493e-16
That is the whole point of the selection: the frame now spans the gate's image.
\end{lstlisting}

\noindent
The two zeros on the losslessness line are the heart of the design. $\Uex \hat S_0 \Vex$
reproduces $\Theb$ exactly and $\|\hat S_0\| = \|\Theb\|$, so nothing has been added to the
state --- only to the space the state is allowed to move in. It also explains why an expansion
alone cannot raise the Schmidt rank, and hence why the evolution has to happen before the
truncation.

The residual line is the discriminating measurement of the whole method: the fraction of the
gate's action on the block that the frame cannot represent, before and after. It falls
\textbf{ten orders of magnitude}, $7.8\times10^{-6} \to 2.5\times10^{-16}$ --- from ``the frame
misses a millionth of the action'' to ``the frame spans it to machine precision''. Driving this
to zero is the strongest thing an expansion at this bond can achieve, and it is exactly the
criterion the two selection defects of Sec.~\ref{sec:retraction} were failing: they left
residuals of $0.19$ and $9.97\times10^{-3}$ untouched \emph{with room to spare}, which is why the
diagnostic is printed as a number rather than asserted as a pass.

Both sides gained four directions ($7 \to 11$) here, and the symmetry is a consequence of the
per-side selection rather than an imposed constraint --- each side spent its own budget against
its own candidates.

\section{Stage 7 --- the $S$-step}

\begin{lstlisting}
S_hat0 (11 x 11)       ->  S1 (11 x 11)        in 14 Lanczos iterations
||S1|| = 1.0000000000   ||S_hat0|| = 1.0000000000    (unitary step: equal)
S1 singular values: 9.950e-01  9.978e-02  5.736e-05  5.752e-06  4.946e-09  4.957e-10  2.852e-13  2.025e-13
kept rank 10 of the expanded 11; discarded weight 2.41e-17
The rank SURVIVES here and only here: S1 is dense on the new directions, so the
truncation keeps them. Without this step they would collapse at the next gauge move.
\end{lstlisting}

\noindent
The norm is preserved to ten digits because the step is a unitary exponential in the expanded
basis. The spectrum of $\hat S_1$ is now \emph{dense on the new directions} --- this is the
moment the expansion turns into rank --- and the truncation therefore keeps them rather than
discarding them as it would have discarded the identically-zero block of $\hat S_0$.

Read the spectrum as the story of the step: two large values ($0.995$, $0.0998$) are the physics
the bond already carried, and the tail from $5.7\times10^{-5}$ down to $2\times10^{-13}$ is what
the expansion just made available. Ten of the eleven directions survive the cutoff, at a
discarded weight of $2.4\times10^{-17}$.

\paragraph{The net effect of one bond update.} Rank $7 \to 11$ by expansion (lossless, no
amplitude changed), then $11 \to 10$ by truncation after the evolution: the bond ends at rank
$\mathbf{10}$, having grown by three. The $14$ Lanczos iterations are the only exponential
anywhere in the update --- there is no $K$-ODE and no $L$-ODE.

\section{Stages 8 and 9 --- the public entry point, and $\tau = 0$}

\begin{lstlisting}
left_core (6 x 2 x 10)     right_core (10 x 2 x 6)    
aug_k = 11  aug_l = 11  keep = 10  n_matvec = 14  discarded = 2.41e-17
err_pre = 5.322e-12  err_fnl = 0.000e+00

||Theta(tau=0) - Theta0|| = 3.95e-15
\end{lstlisting}

\noindent
The entry point reproduces the hand-assembled sequence, and $\tau = 0$ returns the state
unchanged to machine precision --- the sanity check that the expansion has not perturbed
anything, stated as an equation the code must satisfy rather than as an argument.


\end{document}
